% from ega4-2-fr.tex
\oldpage[IV-2]{30}
un morphisme fidèlement plat et quasi-compact (2.2.13 et 1.1.2), on peut se
borner au cas où $S=Y$, $S'=Y'$ en vertu de (I, 3.3.11).

\emph{(i)} Dire que $f$ est séparé signifie que le morphisme diagonal
$\Delta_f:X\to X\times_Y X$ est fermé~; comme
$\Delta_{f'}=(\Delta_f)_{(Y')}$ (I, 5.3.4), si $\Delta_{f'}$ est fermé, il
en est de même de $\Delta_f$ en vertu de (2.6.2), donc $f$ est séparé.

\emph{(ii)} a déjà été démontré sous des hypothèses plus faibles (1.2.5).

\emph{(iii)} et \emph{(iv)}~: La question est évidemment locale sur $X$ et
$Y$, et, compte tenu de (2.2.12), il suffit donc de démontrer le

\begin{lemma}[2.7.1.1]
\label{IV.2.7.1.1-fr}
Soient $A$ un anneau, $B$ une $A$-algèbre, $A'$ une $A$-algèbre qui soit un
$A$-module fidèlement plat, $B'=B\otimes_A A'$. Pour que $B$ soit une
$A$-algèbre de type fini (resp. de présentation finie), il faut et il suffit
que $B'$ soit une $A'$-algèbre de type fini (resp. de présentation finie).
\end{lemma}

On sait déjà que la condition est nécessaire sans hypothèse sur $A'$
(1.3.4, 1.3.6, 1.4.3 et 1.4.6). Supposons que $B'$ soit une $A'$-algèbre de
type fini~; soit $(B_\alpha)_{\alpha\in I}$ la famille filtrante croissante
des sous-$A$-algèbres de $B$, de sorte que l'on a
$B=\varinjlim B_\alpha$, et par suite aussi
$B'=\varinjlim(B_\alpha\otimes_A A')$, le produit tensoriel permutant aux
limites inductives~; si $(x'_i)$ est un système fini de générateurs de la
$A'$-algèbre $B'$, il existe un indice $\alpha$ tel que tous les $x'_i$
appartiennent à la sous-algèbre $B_\alpha\otimes_A A'$ de $B'$, d'où
$B'=B_\alpha\otimes_A A'$, et comme $A'$ est fidèlement plat,
$B=B_\alpha$ $(0_{\mathrm I},6.4.1)$.

Supposons maintenant que $B'$ soit une $A'$-algèbre de présentation finie~;
on sait déjà d'après ce qui précède que $B$ est une $A$-algèbre de type fini,
donc il existe une $A$-algèbre de polynômes
$C=A[T_1,\ldots,T_m]$ et un $A$-homomorphisme surjectif d'algèbres
$C\to B$~; soit $\mathfrak J$ le noyau de cet homomorphisme, de sorte que
l'on a une suite exacte
\[
0\to\mathfrak J\to C\to B\to0,
\]
et par suite aussi une suite exacte $0\to\mathfrak J'\to C'\to B'\to0$
(puisque $A'$ est $A$-plat), en posant
$C'=C\otimes_A A'=A'[T_1,\ldots,T_m]$, et
$\mathfrak J'=\mathfrak J\otimes_A A'$ (identifié à un idéal de $C'$).
Comme $B'$ est une $A'$-algèbre de présentation finie, $\mathfrak J'$ est un
$C'$-module de type fini (1.4.4)~; mais on a
$\mathfrak J'=\mathfrak J\otimes_C C'$ et $C'$ est un $C$-module fidèlement
plat (2.2.13 et 2.2.3)~; on sait alors que l'hypothèse que $\mathfrak J'$ soit
un $C'$-module de type fini entraîne que $\mathfrak J$ est un $C$-module de
type fini (Bourbaki, \emph{Alg. comm.}, chap.~I\textsuperscript{er}, \S~3,
n\textsuperscript{o}~6, prop.~11), donc $B$ est une $A$-algèbre de
présentation finie.

\emph{(v)} résulte de (iii) et de (2.6.2, (v)) en vertu de (1.5.2).

\emph{(vi)} résulte de même de (iv), (v) et (ii) en vertu de (1.6.1).

\emph{(vii)} résulte de (i), (v) et (2.6.4, (ii)) (II, 5.4.1).

\emph{(viii)} Notons d'abord que comme $f'$ est un isomorphisme, c'est un
homéomorphisme universel, donc il en est de même de $f$ (2.6.4)~; on en
conclut déjà que $f$ est quasi-compact et séparé (2.4.4). Posons
$f=(\psi,\theta)$, où $\psi$ est donc un homéomorphisme~; il s'agit de
prouver que
\[
\theta:\mathcal O_Y\longrightarrow f_*(\mathcal O_X)
\]
est un isomorphisme de $\mathcal O_Y$-Modules. Or, si l'on pose
$f'=(\psi',\theta')$, l'homomorphisme
$\theta':\mathcal O_{Y'}\to f'_*(\mathcal O_{X'})$ est composé de
l'homomorphisme canonique
$g^*(f_*(\mathcal O_X))\to f'_*(\mathcal O_{X'})$ et de $g^*(\theta)$
(2.3.2)~; mais le premier de ces deux homomorphismes est bijectif en vertu de
l'hypothèse sur $g$ (2.3.1), donc si $\theta'$ est bijectif,
